\chapter{Introduction and Method}
\section{Problem and purpose}
An online clothing purchase commits to a specific garment variant, not merely a product photograph. A shirt design can have several colour and size combinations, each with its own stock availability. An information system that records only the design cannot reliably decide whether a selected variant can be sold, picked, or exchanged. The purpose of this report is to identify and analyse the requirements for an integrated fashion ecommerce system and translate those requirements into a coherent design.

The proposed system is called \emph{Atelier Fashion Ecommerce}. Atelier Fashion is a fictional retailer that owns its merchandise and sells directly to customers. Its commercial objective is to support dependable purchases and after-sales service while giving staff a consistent view of catalogue, inventory, orders, and financial events. The study emphasises correctness at the boundaries between business processes: checkout and payment, payment and packing, and return inspection and stock restoration.

\section{Evidence and limitations}
Three student reports supplied in \texttt{third\_party} were inspected: the fashion investigation by L\^e Nh\u{u} Qu\`ynh, a marketplace report by Chu Tuyet Nhi, and a retail report by Ph\`ung Thu H\u{o}\u{o}ng. Their shared sections establish a useful coverage baseline: business investigation, stakeholder needs, information analysis, problems, proposed system, requirements, use cases, activity and swimlane diagrams, traceability, and validation. The marketplace's independent sellers and payouts and the other retail report's expiry-oriented inventory are not transferred to this single-retailer fashion case. Neither example is treated as proof of UML correctness or a binding assessment rubric.

No interviews, warehouse observations, production measurements, or institution-specific marking rubric were provided. Accordingly, the as-is situation is a plausible scenario for analysis; every numerical target and policy below is a proposed assumption requiring stakeholder approval before implementation. External documentation supports modelling principles and integration patterns, not claims that a real organisation already uses the proposed design. This distinction prevents imagined survey findings or achieved benefits from being presented as evidence.

\section{Analytical approach and theoretical foundation}
Requirements identification establishes the business problem, system boundary, stakeholder goals, and information that must be available. Requirements analysis resolves ambiguities and expresses observable behaviour through use cases, workflows, domain concepts, and rules. System design then introduces technical responsibilities, persistence structures, interfaces, and transaction mechanisms. A database table is therefore not automatically an analysis class, and an internal service is not automatically an external actor.

UML 2.5.1 is the normative reference for UML notation\cite{uml}. This report applies its concepts as follows. Actors represent roles outside the software boundary. Use cases describe goals that yield observable value. An \texttt{include} relationship points from the including use case to mandatory reused behaviour; it does not express chronological order. Actor generalisation means the specialised role can participate in the general role's interactions. Activities model control and responsibility; mutually exclusive decision branches carry guards, and alternative paths rejoin at merge nodes. State machines describe valid event-driven changes for a particular object. Class associations state domain links and multiplicities rather than arbitrary data flow.

Shopify's public model illustrates the practical distinction between a parent product and a size/colour variant, with variant-level commercial information and inventory links\cite{shopifyvariant,shopifyinventory}. This supports the choice to make SKU identity central to Atelier Fashion. It does not require adopting Shopify or copying its complete data schema. Stripe's documented payment and event behaviours motivate separate payment status, callback verification, and retry handling\cite{stripeintent,stripewebhook,stripeidempotent}. A production payment provider must still be selected for the retailer's operating market.

\section{Scope assumptions}
The first release supports one retailer, one warehouse, VND prices, and one shipment per order. Visitors may browse; authenticated customers may maintain a cart, check out, inspect their own orders, cancel eligible orders, and request returns or exchanges. Online prepaid payment is supported through a hosted provider interface. Staff maintain catalogue data, receive stock, pack orders, inspect returns, reconcile refunds, and review reports. Size charts are provided as reference information; automated body-measurement prediction is outside scope.

The retailer proposes a fourteen-day return request window from confirmed delivery, subject to item condition and disclosed product exclusions. An exchange is limited to an alternative size or colour of the same product with the same agreed item price. Exchanges to a different product or price are handled as a return/refund followed by a new purchase. These are case policies, not asserted universal ecommerce rules or legal conclusions. Guest checkout, cash on delivery, seller onboarding, seller settlement, international tax calculation, manufacturing, and procurement accounting are outside the first release.
